\section{Trace-pairing expansion in transfer-matrix form}
\label{sec:trace_pairing_expansion}
%% =========================================================

% Formula/source: Wolf, Eq. (2.20), writes the transfer coefficients for
% Hilbert--Schmidt orthonormal bases as
% \widehat T_{\alpha\beta}=\tr(F_\alpha^\dagger T(G_\beta)).  Identifying the
% two bases with a Hermitian Hilbert--Schmidt orthonormal basis {\sigma_i}
% gives t_{ij}=\tr(\sigma_i T(\sigma_j)).  The chapter proves the analogous
% coordinate expansion more generally for a basis self-dual under the
% bilinear pairing (A,B)\mapsto\tr(AB):
% T(\rho)=\sum_{i,j}t_{ij}\tr(\sigma_j\rho)\sigma_i.
% Index verdict: j is summed in the input coordinate \tr(\sigma_j\rho), while
% i is summed in the coefficient and output basis element \sigma_i.  This is an
% algebraic basis expansion, not one tensor-network contraction.  In the
% retired figure only \sigma_j and \rho were contracted; t_{ij}, \sigma_i, and
% the equality to T(\rho) were unattached labels.

\begin{definition}[Trace-self-dual basis]\label{def:trace_pairing_self_dual_basis}
    \lean{TracePairingSelfDualBasis}
    \leanok
    A basis $\{\sigma_i\}_i$ of $\MN{D}$ is \emph{trace-self-dual} when
    its coordinate functionals are given by trace pairing:
    \begin{align}
        X &= \sum_i \tr(\sigma_i X)\sigma_i.
        \label{eq:representations_self_dual}
    \end{align}
\end{definition}

\begin{definition}[Trace-pairing coefficients]\label{def:trace_pairing_coeff}
    \lean{tracePairingCoeff}
    \leanok
    \uses{def:trace_pairing_self_dual_basis}
    For a linear map $T : \MN{D} \to \MN{D}$ and a trace-self-dual basis
    $\{\sigma_i\}$, define coefficients
    $t_{ij}:=\tr(\sigma_i T(\sigma_j))$.
\end{definition}

\begin{theorem}[Trace-pairing expansion of a linear map]\label{thm:linearMap_expand_tracePairing}
    \lean{linearMap_expand_tracePairing}
    \leanok
    \uses{def:trace_pairing_self_dual_basis, def:trace_pairing_coeff}
    If $\{\sigma_i\}$ is trace-self-dual, then every linear map
    $T : \MN{D} \to \MN{D}$ admits the expansion
    \begin{align}
        T(\rho)
        &= \sum_i \sum_j t_{ij}\tr(\sigma_j\rho)\sigma_i,
        \label{eq:representations_trace_expansion}\\
        t_{ij}
        &= \tr(\sigma_i T(\sigma_j)).
        \notag
    \end{align}
\end{theorem}

\begin{proof}
    \leanok
    Expand $\rho$ and each $T(\sigma_j)$ in the chosen basis and substitute.
    The trace-self-dual identity (\ref{eq:representations_self_dual}) replaces
    coordinates by traces, giving the double-sum formula
    (\ref{eq:representations_trace_expansion}).
\end{proof}

%% =========================================================
\section{General matrix singular value decompositions}
\label{sec:svd_normal_form}
%% =========================================================

The results in this section are ordinary singular value decompositions in the
general complex matrix-algebraic setting. They are useful for transfer matrices, but
they are not the real $3\times3$ singular value decomposition arising from the
    $\mathrm{SO}(3)$ action on the traceless Pauli block of a qubit channel in
\cite[Section~2.4]{Wolf2012Quantum}. The Lorentz normal form theorem
(Theorem~\ref{thm:lorentz_normal_form_qubit}) instead uses general invertible
Kraus-rank-one completely positive filterings. Its proof requires the
minimisation argument and Lorentz-orbit analysis in
\cite[Propositions~2.9 and~2.11]{Wolf2012Quantum}; see
Section~\ref{sec:lorentz_normal_form}.

\begin{theorem}[SVD for positive semidefinite matrices]
    \label{thm:svd_posSemidef}
    \lean{Matrix.svd_of_posSemidef}\leanok
    Every positive semidefinite matrix $M \in \MN{D}$ admits a decomposition
    $M=U\diag(\sigma)U^{\dagger}$, where $U \in \mc{U}(D)$ is unitary and
    $\sigma : \{0,\ldots,D-1\} \to \R_{\ge 0}$ has non-negative entries.
\end{theorem}

\begin{proof}\leanok
    The spectral theorem for Hermitian matrices provides an orthonormal
    eigenbasis with real eigenvalues. Positive semidefiniteness forces those
    eigenvalues to be non-negative, and the decomposition is written in SVD
    form with $U$ the eigenvector unitary and $\sigma$ the eigenvalue sequence.
\end{proof}

\begin{theorem}[SVD existence for invertible matrices]
    \label{thm:svd_of_isUnit}
    \lean{Matrix.svd_of_isUnit}\leanok
    Every invertible complex square matrix $M \in \MN{D}$ admits a singular
    value decomposition $M=U\diag(\sigma)V^{\dagger}$, with
    $U,V \in \mc{U}(D)$ unitary and $\sigma_i>0$ for all $i$.
\end{theorem}

\begin{proof}\leanok
    Apply the spectral theorem to the positive definite matrix
    $M^{\dagger}M$ to obtain
    $M^{\dagger}M=V\diag(\lambda)V^{\dagger}$ with $\lambda_i>0$.
    Set $\sigma_i:=\sqrt{\lambda_i}$ and $\Sigma:=\diag(\sigma)$; since
    $\Sigma$ is a real positive diagonal, it is self-adjoint and invertible.
    Define $U:=MV\Sigma^{-1}$. Then
    \begin{align}
        U^{\dagger}U
        &= \Sigma^{-1}V^{\dagger}(M^{\dagger}M)V\Sigma^{-1}
         = \Sigma^{-1}\diag(\lambda)\Sigma^{-1}
         = \Id,
        \notag\\
        U\Sigma V^{\dagger}
        &= MV(\Sigma^{-1}\Sigma)V^{\dagger}
         = MVV^{\dagger}
         = M.
        \notag
    \end{align}
    Thus $U$ is unitary and the stated decomposition holds.
\end{proof}

\begin{theorem}[SVD representation of a transfer matrix]
    \label{thm:transferMatrix_svd_of_isUnit}
    \lean{transferMatrix_svd_of_isUnit}\leanok
    \uses{thm:svd_of_isUnit}
    Every invertible transfer matrix $\widehat{T}$ of a linear super-operator
    $T : \MN{D} \to \MN{D}$ admits a singular value decomposition
    $\widehat{T}=U\diag(\sigma)V^{\dagger}$ with $U,V$
    unitary on $\C^{D}\otimes\C^{D}$ and $\sigma_i>0$.
\end{theorem}

\begin{proof}\leanok
    Direct specialisation of Theorem~\ref{thm:svd_of_isUnit} to the index type
    $\{0,\ldots,D-1\}\times\{0,\ldots,D-1\}$ used by the transfer-matrix
    representation.
\end{proof}

\begin{remark}[Sorted / unique singular values]
    Theorem~\ref{thm:svd_of_isUnit} produces the singular values as an
    unordered family, without the usual convention that $\sigma_i$ are sorted
    in non-increasing order or the statement that they are uniquely determined
    by $M$. Applications that distinguish individual singular values may also
    require a non-increasing ordering and uniqueness of the resulting ordered
    family. These refinements are not part of the present statement.
\end{remark}

%% =========================================================
\section{Lorentz normal form}
\label{sec:lorentz_normal_form}
%% =========================================================

This section states the existence theorems from
\cite[Section~2.4]{Wolf2012Quantum}.
% The existing Wolf.* declarations referenced in this section live in
% QICLean/Channel/LorentzNormalForm.lean. The correctly formulated Proposition 2.11
% has no Lean declaration yet; declaration tags are omitted for pending results.

\subsection{Canonical metric blocks for the Minkowski reduction}

Let $P_n$ be the reversal permutation matrix, so that
$(P_n)_{ij}=1$ exactly when $i+j=n+1$.  For a sign
$\varepsilon\in\{1,-1\}$, the matrix $\varepsilon P_n$ is the signed
metric block attached to a real-eigenvalue Jordan block in the sign
characteristic of Gohberg--Lancaster--Rodman, Theorem~5.3.

\begin{definition}[Signed sip metric blocks]
    \lean{Matrix.GLRSign, Matrix.sipMatrix, Matrix.signedSipMatrix}
    \leanok
    The \emph{sip matrix} of size $n$ is $P_n$.  A signed sip block is
    $\varepsilon P_n$, where $\varepsilon=1$ or $\varepsilon=-1$.
    Each such block is symmetric, invertible, and its own inverse.
\end{definition}

\begin{definition}[Real canonical Jordan blocks]
    \lean{Matrix.realJordanBlock, Matrix.realConjugatePairBlock,
        Matrix.realConjugatePairJordanBlock, Matrix.complexPairSipMatrix}
    \leanok
    A real-eigenvalue block is the ordinary real Jordan block $J_n(a)$
    with eigenvalue $a$ and ones on its first superdiagonal.  If
    $a\pm\tau i$ is a non-real conjugate pair, put
    \begin{align}
        B(a,\tau)&=\begin{pmatrix}a&\tau\\-\tau&a\end{pmatrix}.
    \end{align}
    Its real Jordan chain of length $m$ has $B(a,\tau)$ on the block
    diagonal and $I_2$ on the first block superdiagonal.  On the product
    indexing of this $2m$-dimensional real block, its unsigned metric
    reverses both the chain coordinate and the two-dimensional
    realification coordinate.  No sign characteristic is attached to a
    non-real block.
\end{definition}

\begin{theorem}[Canonical blocks are selfadjoint for their metric blocks]
    \lean{Matrix.signedSipMatrix_mul_realJordanBlock,
        Matrix.complexPairSipMatrix_mul_realConjugatePairJordanBlock}
    \leanok
    For every sign $\varepsilon$ and every real Jordan block,
    \begin{align}
        (\varepsilon P_n)J_n(a)
        &=J_n(a)^{\mathsf T}(\varepsilon P_n).
    \end{align}
    The real Jordan chain for a non-real conjugate pair satisfies the
    analogous identity with its unsigned $2m$-dimensional reversal metric.
\end{theorem}

\begin{definition}[Four-dimensional Minkowski selfadjointness]
    \lean{Matrix.minkowskiMetric, Matrix.IsMinkowskiSelfAdjoint}
    \leanok
    Let
    \begin{align}
        M&=\operatorname{diag}(1,-1,-1,-1).
    \end{align}
    A real matrix $C\in M_4(\mathbb R)$ is \emph{$M$-selfadjoint} when
    $C^{\mathsf T}M=MC$.
\end{definition}

\begin{definition}[The four VDDM block patterns]
    \lean{Matrix.GLRFourDimensionalPattern}
    \leanok
    The four source-indexed patterns are: four real one-dimensional
    blocks; one two-dimensional non-real conjugate-pair block and two real
    one-dimensional blocks; one real Jordan block of size two and two real
    one-dimensional blocks; or one real Jordan block of size three and one
    real one-dimensional block.
\end{definition}

The load-bearing existence statement is still open.  It must construct, for
every $M$-selfadjoint $C$, an invertible $X$, a real canonical Jordan matrix
$J$, and its signed metric $N_J$ satisfying
\begin{align}
    C&=X^{-1}JX,
    &XMX^{\mathsf T}&=N_J,
\end{align}
and then derive exhaustiveness of the four patterns from inertia $(1,3)$.
The present library has neither real Jordan-form existence nor the simultaneous
similarity--congruence theorem controlling the sign characteristic.  The exact
formalization boundary and the required transpose and inverse translations are
recorded in
\path{docs/paper-gaps/qiclean_minkowski_canonical_pair_dependency.tex}.

\begin{definition}[Pauli transfer matrix]\label{def:pauli_transfer_matrix}
    \lean{Wolf.pauliMatrices, Wolf.pauliTransferEntry,
        Wolf.pauliTransferMatrix, Wolf.pauliTransferMatrixReal}
    \leanok
    Let $\sigma_0=\Id,\sigma_1,\sigma_2,\sigma_3$ be the Pauli matrices.  The
    \emph{Pauli transfer matrix} of a linear map
    $T:\MN{2}\to\MN{2}$ is $\widehat T\in M_4(\C)$ with entries
    \begin{align}
        \widehat T_{ij}
        &=\frac12\tr\!\left(\sigma_iT(\sigma_j)\right),
        \qquad 0\leq i,j\leq3.
        \label{eq:pauli_transfer_entries}
    \end{align}
    We write
    \begin{align}
        \widehat T
        &=
        \begin{pmatrix}
            \widehat T_{00} & r^{\mathsf T}\\
            v & \Delta
        \end{pmatrix},
        \label{eq:pauli_transfer_blocks}
    \end{align}
    where $\Delta=(\widehat T_{ij})_{i,j=1}^{3}$.  Thus $r^{\mathsf T}$ and
    $v$ are exactly the two off-diagonal Pauli blocks.  This is the
    representation used in \cite[Section~2.4, Eq.~(2.39)]{Wolf2012Quantum}.
\end{definition}

\begin{definition}[Pauli-block diagonal truncation]
    \label{def:pauli_block_diagonal_truncation}
    \lean{Wolf.pauliTimeReversal, Wolf.pauliTimeReversalDiagonal,
        Wolf.pauliBlockDiagonalTruncation}
    \leanok
    \uses{def:pauli_transfer_matrix}
    Define Pauli time reversal by
    $\Theta(X)=\tr(X)\Id-X$.  Its Pauli transfer matrix is
    \begin{align}
        D&=\diag(1,-1,-1,-1).
        \label{eq:pauli_time_reversal_diagonal}
    \end{align}
    For a linear map $T:\MN{2}\to\MN{2}$, define its
    \emph{Pauli-block diagonal truncation} by
    \begin{align}
        T'&=\frac{T+\Theta\circ T\circ\Theta}{2}.
        \label{eq:pauli_block_truncation}
    \end{align}
    This is the truncation in
    \cite[Proposition~2.10, Section~2.4]{Wolf2012Quantum}; see also the local
    source \path{Notes/WolfNoteTexSource/ch02_representations.tex},
    lines~984--998.
\end{definition}

\begin{theorem}[Pauli-block truncation preserves positivity and complete positivity]
    \label{thm:pauli_block_truncation_forward}
    \lean{Wolf.pauliTransferEntry_timeReversal_comp,
        Wolf.pauliTransferEntry_pauliBlockDiagonalTruncation,
        Wolf.pauliBlockDiagonalTruncation_isPositiveMap,
        Wolf.pauliBlockDiagonalTruncation_isCPMap}
    \leanok
    \uses{def:pauli_block_diagonal_truncation, def:positive_map, def:cp_map}
    Let $T:\MN{2}\to\MN{2}$ be Hermiticity preserving, and let $T'$ be its
    Pauli-block diagonal truncation.  If $\widehat T$ is written as in
    \eqref{eq:pauli_transfer_blocks}, then
    \begin{align}
        \widehat{T'}
        &=\frac{\widehat T+D\widehat T D}{2}
          =\begin{pmatrix}
              \widehat T_{00}&0\\
              0&\Delta
            \end{pmatrix}
          =\widehat T_{00}\oplus\Delta.
        \label{eq:pauli_block_truncation_matrix}
    \end{align}
    Equivalently, entrywise,
    \begin{align}
        \widehat{T'}_{ij}
        &=\begin{cases}
            \widehat T_{ij},& i=j=0\text{ or }i,j\in\{1,2,3\},\\
            0,&\text{otherwise}.
          \end{cases}
        \label{eq:pauli_block_truncation_entrywise}
    \end{align}
    If $T$ is positive, then $T'$ is positive.  If $T$ is completely
    positive, then $T'$ is completely positive.  These are precisely the
    forward implications proved in
    \cite[Proposition~2.10, Section~2.4]{Wolf2012Quantum}, following
    Eq.~(2.39); the local proof is at
    \path{Notes/WolfNoteTexSource/ch02_representations.tex}, lines~992--998.
\end{theorem}

\begin{proof}\leanok
    In the Pauli basis, pre- and postcomposition by $\Theta$ multiply the
    transfer matrix by $D$ on the right and left, respectively.  Averaging
    with $D\widehat T D$ therefore gives
    \eqref{eq:pauli_block_truncation_entrywise}.  Time reversal is positive,
    so $\Theta\circ T\circ\Theta$ is positive whenever $T$ is positive, and
    convexity gives positivity of $T'$.  If
    $T(X)=\sum_iK_iXK_i^\dagger$, then
    $\Theta\circ T\circ\Theta$ has Kraus operators
    $\sigma_2\overline{K_i}\sigma_2$.  Hence the same averaging also
    preserves complete positivity.
\end{proof}

\begin{remark}[Source correction: the converse fails]
    \label{rem:pauli_block_truncation_converse}
    Proposition~2.10 of \cite{Wolf2012Quantum} prints both preservation
    statements as equivalences, but its proof establishes only the forward
    implications in Theorem~\ref{thm:pauli_block_truncation_forward}.  Indeed,
    consider
    \begin{align}
        T(X)&=\tr(X)\diag(3/2,-1/2).
        \label{eq:pauli_block_truncation_counterexample}
    \end{align}
    This map is Hermiticity preserving but not positive, since
    $T(\Id/2)=\diag(3/2,-1/2)$.  Its Pauli-block diagonal truncation is the
    completely depolarizing channel
    $T'(X)=\tr(X)\Id/2$, which is completely positive.  Thus neither the
    positivity converse nor the complete-positivity converse is valid.  See
    \path{docs/paper-gaps/wolf_prop2_10_block_truncation_converse.tex} for the
    focused source-correction record.
\end{remark}

\begin{definition}[SL-filtering operation]\label{def:sl_filtering}
    \lean{Wolf.SLFiltering}\leanok
    An \emph{SL-filtering} for $D \times D$ matrices is a completely
    positive map of the form $\Phi(X)=SXS^{\dagger}$, where
    $S \in \MN{D}$ satisfies $\det S=1$. Such maps are invertible; when
    $D\geq1$, they have Kraus rank~1.
\end{definition}

\begin{definition}[Invertible filtering operation]
\label{def:invertible_filtering}
    \lean{Wolf.InvertibleFilter,
        Wolf.InvertibleFilter.map,
        Wolf.InvertibleFilter.map_apply,
        Wolf.InvertibleFilter.map_eq_unitaryConjLM,
        Wolf.InvertibleFilter.cp,
        Wolf.InvertibleFilter.id,
        Wolf.InvertibleFilter.map_id,
        Wolf.InvertibleFilter.comp,
        Wolf.InvertibleFilter.comp_X,
        Wolf.InvertibleFilter.map_comp,
        Wolf.InvertibleFilter.inv,
        Wolf.InvertibleFilter.inv_X,
        Wolf.InvertibleFilter.inv_comp,
        Wolf.InvertibleFilter.comp_inv,
        Wolf.InvertibleFilter.inv_map_comp,
        Wolf.InvertibleFilter.map_comp_inv,
        Wolf.InvertibleFilter.filteredMap,
        Wolf.InvertibleFilter.filteredMap_apply,
        Wolf.SLFiltering.toInvertibleFilter,
        Wolf.SLFiltering.toInvertibleFilter_X,
        Wolf.SLFiltering.toInvertibleFilter_map}\leanok
    An \emph{invertible filtering operation} on $\MN{D}$ is a completely
    positive map
    \begin{align}
        \Phi_X(A)=XAX^\dagger,
        \qquad X\in\GL(D,\C).
        \label{eq:representations_invertible_filter}
    \end{align}
    If $T:\MN{d_1}\to\MN{d_2}$, pre- and postfiltering are written in
    Wolf's order as $\Phi_2\circ T\circ\Phi_1$.  The matrices $X$ are
    retained as part of the data; in particular, matrices which differ by a
    phase are not identified.
\end{definition}

\begin{lemma}[Scalar and determinant-one decomposition]
\label{lem:invertible_filter_scalar_sl}
    \lean{Wolf.InvertibleFilter.hasKrausCard_one,
        Wolf.InvertibleFilter.map_ne_zero,
        Wolf.InvertibleFilter.choiRank_eq_one,
        Wolf.InvertibleFilter.singleKrausMap_smul,
        Wolf.InvertibleFilter.exists_scalar_slFiltering}\leanok
    Let $D\geq1$ and let $X\in\GL(D,\C)$.  There are a nonzero
    $c\in\C$ and $S\in\SL(D,\C)$ such that
    \begin{align}
        c^D&=\det X, & X&=cS,
        \label{eq:representations_filter_scalar_matrix}\\
        \Phi_X&=|c|^2\Phi_S, & |c|^2&>0.
        \label{eq:representations_filter_scalar_map}
    \end{align}
    Moreover, $\Phi_X$ has Kraus rank exactly one.  The matrix scalar $c$
    is complex and need not be positive real; positivity applies only to the
    map scalar $|c|^2$.
\end{lemma}

\begin{proof}\leanok
    Choose a $D$th root $c$ of $\det X$.  Since $X$ is invertible, both
    $\det X$ and $c$ are nonzero.  Put $S=c^{-1}X$.  Multiplicativity of the
    determinant gives $\det S=1$, while direct expansion of
    $(cS)A(cS)^\dagger$ gives
    $\Phi_X=|c|^2\Phi_S$.  The displayed nonzero Kraus operator gives Kraus
    rank at most one, and invertibility excludes Kraus rank zero.
\end{proof}

\begin{lemma}[Scalar factors under pre- and postfiltering]
\label{lem:invertible_filter_two_scalar}
    \lean{Wolf.InvertibleFilter.filteredMap_eq_normSq_smul}\leanok
    \uses{def:invertible_filtering, lem:invertible_filter_scalar_sl}
    Suppose $X_i=c_iS_i$, where $c_i\neq0$ and
    $S_i\in\SL(d_i,\C)$.  For every linear map
    $T:\MN{d_1}\to\MN{d_2}$,
    \begin{align}
        \Phi_{X_2}\circ T\circ\Phi_{X_1}
        =|c_1c_2|^2
          \bigl(\Phi_{S_2}\circ T\circ\Phi_{S_1}\bigr).
        \label{eq:representations_filter_two_scalar}
    \end{align}
\end{lemma}

\begin{proof}\leanok
    Substitute $\Phi_{X_i}=|c_i|^2\Phi_{S_i}$ and use linearity of $T$ and
    of the outer filtering operation.  The two real scalars combine as
    $|c_1|^2|c_2|^2=|c_1c_2|^2$.
\end{proof}

\begin{definition}[Doubly-stochastic map]\label{def:doubly_stochastic}
    \lean{Wolf.DoublyStochastic}\leanok
    A linear map $T : \MN{D} \to \MN{D}$ is \emph{doubly-stochastic}
    if $T(\Id)\propto\Id$ and the reduced density matrix
    $\tr_{1}[\tau]$ of its Choi matrix
    $\tau=(T\otimes\id)(\ketbra{\Omega}{\Omega})$
    is proportional to the identity. This is the normal form in
    \cite[Proposition~2.9]{Wolf2012Quantum}.
\end{definition}

\begin{lemma}[Positive-definite trace lower bound]\label{lem:posdef_trace_lower_bound}
    \lean{posDef_minEigenvalue_mul_trace_conjTranspose_mul_self_le,
        sub_minEigenvalue_smul_one_posSemidef}
    \leanok
    Let $D\geq1$, let $M\in\MN{D}$ be positive-definite, and let
    $\lambda_{\min}(M)$ be its smallest Hermitian eigenvalue. For every
    $X\in\MN{D}$,
    \begin{align}
        \lambda_{\min}(M)\tr(X^{\dagger}X)
        &\leq \tr(XMX^{\dagger}).
        \label{eq:representations_posdef_trace_bound}
    \end{align}
\end{lemma}

\begin{proof}
    \leanok
    The matrix $M-\lambda_{\min}(M)\Id$ is positive semidefinite by the
    spectral theorem. Since $X^{\dagger}X$ is also positive semidefinite, the
    trace product
    $\tr\!\left(X^{\dagger}X(M-\lambda_{\min}(M)\Id)\right)$
    is non-negative. Expanding this expression and cycling the trace gives
    (\ref{eq:representations_posdef_trace_bound}).
\end{proof}

\begin{lemma}[Determinant AM--GM Hilbert--Schmidt bound]
\label{lem:det_amgm_hilbert_schmidt_bound}
    \lean{Matrix.card_le_trace_conjTranspose_mul_self_re_of_det_norm_eq_one}
    \leanok
    Let $A\in\MN{D}$ with $D\geq1$. If $|\det A|=1$, then
    $D\leq\tr(A^{\dagger}A)$.
\end{lemma}

\begin{proof}
    \leanok
    The matrix $A^{\dagger}A$ is positive semidefinite and has determinant
    $|\det A|^{2}=1$. Thus the product of its Hermitian eigenvalues is one.
    Denoting the eigenvalues of $A^{\dagger}A$ by
    $\lambda_{1},\ldots,\lambda_{D}\geq0$, the arithmetic-geometric mean
    inequality gives
    \begin{align}
        \frac{\lambda_{1}+\cdots+\lambda_{D}}{D}
        &\geq(\lambda_{1}\cdots\lambda_{D})^{1/D}
         =1.
        \notag
    \end{align}
    Therefore
    $\tr(A^{\dagger}A)=\lambda_{1}+\cdots+\lambda_{D}\geq D$.
\end{proof}

\begin{lemma}[Kronecker Hilbert--Schmidt trace multiplicativity]
\label{lem:kronecker_hilbert_schmidt_trace}
    \lean{Matrix.trace_conjTranspose_mul_self_kronecker}
    \leanok
    For complex matrices $A$ and $B$ of compatible rectangular sizes,
    \begin{align}
        \tr\!\left((A\otimes_{k}B)^{\dagger}
            (A\otimes_{k}B)\right)
        &=\tr(A^{\dagger}A)\tr(B^{\dagger}B).
        \label{eq:representations_kronecker_hs}
    \end{align}
\end{lemma}

\begin{proof}
    \leanok
    Use
    $(A\otimes_{k}B)^{\dagger}=A^{\dagger}\otimes_{k}B^{\dagger}$, the mixed
    product identity for Kronecker products, and
    $\tr(C\otimes_{k}D)=\tr(C)\tr(D)$. With
    $C=A^{\dagger}A$ and $D=B^{\dagger}B$, these identities give
    (\ref{eq:representations_kronecker_hs}).
\end{proof}

% Lean: Wolf.infimum_is_attained
\begin{lemma}[Infimum attainment for SL-filterings]\label{lem:infimum_is_attained}
    \lean{Wolf.infimum_is_attained}\leanok
    \uses{lem:posdef_trace_lower_bound, lem:det_amgm_hilbert_schmidt_bound,
    lem:kronecker_hilbert_schmidt_trace}
    For a positive-definite operator $\tau$ on $\C^{d_2}\otimes\C^{d_1}$,
    the infimum of
    \begin{align}
        \tr\!\left[(S_{2}\otimes_{k}S_{1})\tau
            (S_{2}\otimes_{k}S_{1})^{\dagger}\right]
        \notag
    \end{align}
    over $S_{1}\in\MN{d_1}$ and $S_{2}\in\MN{d_2}$ with
    $\det S_{1}=\det S_{2}=1$ is attained.  The two tensor factors may have
    different dimensions.
\end{lemma}

\begin{proof}\leanok
    \uses{lem:posdef_trace_lower_bound, lem:det_amgm_hilbert_schmidt_bound,
    lem:kronecker_hilbert_schmidt_trace}
    Write $X=S_{2}\otimes_{k}S_{1}$.
    Lemma~\ref{lem:posdef_trace_lower_bound} gives the trace lower bound below.
    Lemma~\ref{lem:kronecker_hilbert_schmidt_trace} factors the
    Hilbert--Schmidt trace. For each $i\in\{1,2\}$,
    Lemma~\ref{lem:det_amgm_hilbert_schmidt_bound} supplies the first row.
    Together,
    \begin{align}
        \tr(S_i^\dagger S_i)
        &=\|S_i\|^2\geq d_i,
        \notag\\
        \tr(X\tau X^\dagger)
        &\geq\lambda_{\min}(\tau)\tr(X^\dagger X)
        \notag\\
        &=\lambda_{\min}(\tau)
          \tr(S_2^\dagger S_2)\tr(S_1^\dagger S_1)
        \notag\\
        &\geq\lambda_{\min}(\tau)\,d_j\,\|S_i\|^2
          \qquad (j\neq i).
        \notag
    \end{align}
    Hence any point
    whose value is at most the value at the identity lies in a fixed
    Frobenius ball $\{\,\|S\|\leq C\,\}$. Intersecting with the closed set
    $\det S=1$ produces a compact set, on which the continuous trace
    functional attains its minimum by the extreme-value theorem. A point
    outside the sublevel set has value larger than the value at the identity,
    so the minimiser on the compact set is a global minimiser.
\end{proof}

% Lean: pow_card_mul_prod_le_sum_pow, pow_card_mul_prod_le_sum_pow'
\begin{lemma}[Arithmetic--geometric-mean inequality, product/sum form]
\label{lem:amgm_prod_sum}
    \lean{pow_card_mul_prod_le_sum_pow, pow_card_mul_prod_le_sum_pow'}
    \leanok
    For a non-negative family $f_{0},\ldots,f_{D-1}$ of real numbers, or more
    generally for a non-negative family indexed by a finite set with $D$
    elements,
    \begin{align}
        D^{D}\prod_{i}f_{i}
        &\leq\left(\sum_{i}f_{i}\right)^{D}.
        \label{eq:representations_amgm_prod_sum}
    \end{align}
\end{lemma}

\begin{proof}\leanok
    Apply the weighted arithmetic--geometric-mean inequality with uniform
    weights $1/D$:
    \begin{align}
        \left(\prod_{i}f_{i}\right)^{1/D}
        &\leq\frac{1}{D}\sum_{i}f_{i}.
        \notag
    \end{align}
    Raising both sides to the $D$-th power gives
    (\ref{eq:representations_amgm_prod_sum}).
\end{proof}

% Lean: Matrix.PosSemidef.pow_card_mul_det_re_le_trace_re_pow,
% Matrix.PosSemidef.pow_card_mul_det_re_eq_trace_re_pow_iff,
% Matrix.PosSemidef.pow_card_mul_det_re_le_trace_re_pow'
\begin{lemma}[Trace--determinant AM--GM for positive-semidefinite matrices]
\label{lem:trace_det_amgm}
    \lean{Matrix.PosSemidef.pow_card_mul_det_re_le_trace_re_pow,
    Matrix.PosSemidef.pow_card_mul_det_re_eq_trace_re_pow_iff,
    Matrix.PosSemidef.pow_card_mul_det_re_le_trace_re_pow'}
    \leanok
    \uses{lem:amgm_prod_sum}
    For a positive-semidefinite $D \times D$ matrix $M$,
    \begin{align}
        D^{D}\det M
        &\leq(\tr M)^{D},
        \label{eq:representations_trace_det_amgm}
    \end{align}
    and equality holds if and only if $M=(\tr M/D)\Id$. The inequality holds
    for a positive-semidefinite matrix whose rows and columns are indexed by
    any finite set with $D$ elements; the equality characterization is stated
    only for $D \times D$ matrices.
\end{lemma}

\begin{proof}\leanok
    Both $\det M=\prod_{i}\lambda_{i}$ and
    $\tr M=\sum_{i}\lambda_{i}$ are expressed through the non-negative
    eigenvalues $\lambda_i$ of $M$. The bound is the
    arithmetic--geometric-mean inequality with uniform weights $1/D$, raised
    to the $D$-th power; for a $D \times D$ matrix, equality holds exactly
    when all eigenvalues coincide, that is, when $M$ is a scalar matrix.
    Relabelling along a bijection with the index set of a $D \times D$ matrix
    carries the inequality to an arbitrary finite index set of $D$ elements.
\end{proof}

% Lean: Wolf.exists_normal_form_generic_tau
\begin{theorem}[Normal form for generic $\tau$]
    \label{thm:exists_normal_form_generic_tau}
    \lean{Wolf.exists_normal_form_generic_tau}
    \leanok
    \uses{lem:infimum_is_attained, lem:trace_det_amgm}
    Let $\tau\in\mathcal{B}(\C^{d_2}\otimes\C^{d_1})$ be positive-definite. Then
    there exist $S_i\in\mathrm{SL}(d_i,\C)$ which attain the infimum
    \begin{align}
        p
        &:={\inf}_{X_i\in\mathrm{SL}(d_i,\C)}
          \tr\!\left[(X_2\otimes X_1)\tau
          (X_2\otimes X_1)^\dagger\right].
        \label{eq:representations_generic_tau_infimum}
    \end{align}
    In particular, for
    \begin{align}
        \tau'
        &:=(S_2\otimes S_1)\tau(S_2\otimes S_1)^\dagger,
        \notag
    \end{align}
    both partial traces are proportional to the respective identity matrices:
    \begin{align}
        \tr_{2}[\tau']&=\kappa_1\Id_{d_1},
        &\tr_{1}[\tau']&=\kappa_2\Id_{d_2}
        \label{eq:representations_generic_tau_marginals}
    \end{align}
    for some $\kappa_1,\kappa_2\in\C$. This is
    \cite[Proposition~2.8, lines 894--919]{Wolf2012Quantum}, with the two
    dimensions kept independent.
\end{theorem}

\begin{proof}\leanok
    \uses{lem:infimum_is_attained, lem:trace_det_amgm}
    Choose a minimizing pair $(S_1,S_2)$ by
    Lemma~\ref{lem:infimum_is_attained}. Hold $S_2$ fixed and set
    \begin{align}
        \rho_2
        &:=(S_2\otimes\Id_{d_1})\tau
          (S_2\otimes\Id_{d_1})^\dagger,
        &M_1&:=\tr_2[\rho_2].
        \notag
    \end{align}
    Positive-definiteness of $\tau$ and invertibility of $S_2$ imply that
    $M_1$ is positive-definite. The defining minimality of $(S_1,S_2)$ says
    that $S_1$ minimizes $\tr(XM_1X^\dagger)$ over
    $X\in\mathrm{SL}(d_1,\C)$. A determinant-one whitening of $M_1$ attains
    the trace--determinant arithmetic--geometric-mean bound, while its equality
    case in Lemma~\ref{lem:trace_det_amgm} forces
    $S_1M_1S_1^\dagger$ to be a scalar matrix. This matrix is
    $\tr_2[\tau']$. Holding $S_1$ fixed and interchanging the two tensor
    factors gives $\tr_1[\tau']\propto\Id_{d_2}$ in the same way.
\end{proof}

% Lean: Wolf.exists_normal_form_generic
\begin{theorem}[Generic normal form for CP maps]\label{thm:exists_normal_form_generic}
    \lean{Wolf.exists_normal_form_generic}
    \leanok
    \uses{def:sl_filtering, def:doubly_stochastic, lem:infimum_is_attained,
    lem:trace_det_amgm}
    Let $T : \MN{D} \to \MN{D}$ be a completely positive map whose Choi matrix
    is positive-definite (equivalently, $T$ has full Kraus rank). Then there
    exist SL-filterings $\Phi_{1},\Phi_{2}$ such that
    $\Phi_{2}\circ T\circ\Phi_{1}$ is doubly-stochastic.

    This is the equal-dimension (square) case $T : \MN{D} \to \MN{D}$ of
    \cite[Proposition~2.9]{Wolf2012Quantum}; the rectangular form with
    independent dimensions is
    Theorem~\ref{thm:exists_normal_form_generic_rect}.
    % See docs/paper-gaps/wolf_prop28_square_case.tex (restriction resolved)
\end{theorem}

\begin{proof}\leanok
    \uses{lem:infimum_is_attained, lem:trace_det_amgm}
    Let $\tau$ be the Choi matrix of $T$ and take the minimiser
    $(S_{1},S_{2})$ of
    \begin{align}
        \tr\!\left[(S_{2}\otimes_{k}S_{1})\tau
            (S_{2}\otimes_{k}S_{1})^{\dagger}\right]
        \notag
    \end{align}
    over $\det S_{1}=\det S_{2}=1$ from
    Lemma~\ref{lem:infimum_is_attained}. Set $\Phi_{1},\Phi_{2}$ to be the
    filterings with matrices $S_{1}^{\mathsf T}$ and $S_{2}$, so that the Choi
    matrix of $\Phi_{2}\circ T\circ\Phi_{1}$ equals
    $(S_{2}\otimes_{k}S_{1})\tau
    (S_{2}\otimes_{k}S_{1})^{\dagger}$. Holding one filtering factor fixed,
    the minimiser is optimal in the other coordinate, so each partial trace
    minimises a functional $\tr(SMS^{\dagger})$ over $\det S=1$ for a
    positive-semidefinite $M$, where
    $\tr(SMS^{\dagger})\geq D(\det M)^{1/D}$, with equality at the minimiser;
    by the equality case of Lemma~\ref{lem:trace_det_amgm}, this forces
    $SMS^{\dagger}\propto\Id$, hence both partial traces are proportional to
    the identity, which is exactly the doubly-stochastic condition.
\end{proof}

\begin{definition}[Doubly-stochastic map between matrix algebras]
    \label{def:doubly_stochastic_rect}
    \lean{Wolf.DoublyStochasticRect}\leanok
    A linear map $T : \MN{d_1} \to \MN{d_2}$ between matrix algebras of
    possibly different dimensions is \emph{doubly-stochastic} if
    $T(\Id)\propto\Id$ and $T^{*}(\Id)\propto\Id$, where $T^{*}$ is the
    trace-pairing adjoint.  This is the normal-form condition of
    \cite[Proposition~2.9]{Wolf2012Quantum}; by the rectangular
    Choi--Jamiolkowski correspondence it is equivalent to both partial traces
    of the Choi matrix being proportional to the identity.
\end{definition}

\begin{theorem}[Generic normal form for rectangular CP maps]
    \label{thm:exists_normal_form_generic_rect}
    \lean{Wolf.exists_normal_form_generic_rect}
    \leanok
    \uses{def:sl_filtering, def:doubly_stochastic_rect,
    thm:exists_normal_form_generic_tau, def:choi_matrix_rect, def:is_kraus_cp}
    Let $T : \MN{d_1} \to \MN{d_2}$ be a completely positive map whose
    Choi matrix is positive-definite (equivalently, $T$ has full Kraus
    rank). Then there exist SL-filterings $\Phi_{1}$ on $\MN{d_1}$ and
    $\Phi_{2}$ on $\MN{d_2}$ such that $\Phi_{2}\circ T\circ\Phi_{1}$
    is doubly-stochastic.  This is
    \cite[Proposition~2.9]{Wolf2012Quantum} at the source's generality.
\end{theorem}

\begin{proof}\leanok
    \uses{thm:exists_normal_form_generic_tau}
    Apply Theorem~\ref{thm:exists_normal_form_generic_tau} to the rectangular
    Choi matrix $\tau$ of $T$, and let $S_1,S_2$ be the resulting
    determinant-one matrices. Define $\Phi_1,\Phi_2$ to be the filterings with
    matrices $S_1^{\mathsf T}$ and $S_2$, respectively. The Choi matrix of
    $\Phi_2\circ T\circ\Phi_1$ is then the minimizing representative
    $(S_2\otimes S_1)\tau(S_2\otimes S_1)^\dagger$. Its two partial
    traces are scalar by
    Equation~\eqref{eq:representations_generic_tau_marginals}, and the
    rectangular Choi identities
    $\tr_B(\tau)=T(\Id)/d_1$ and $\tr_A(\tau)=(T^{*}(\Id))^{\mathsf T}/d_1$
    give the doubly-stochastic condition.
\end{proof}

\begin{definition}[Diagonal Lorentz normal form]\label{def:lorentz_diagonal}
    \lean{Wolf.IsLorentzDiagonal}\leanok
    For a completely positive, trace-preserving qubit channel
    $T' : \MN{2} \to \MN{2}$, write $\widehat{T'}$ for its matrix
    in the normalized Pauli basis
    $\{\sigma_{0}/\sqrt{2},\sigma_{1}/\sqrt{2},
        \sigma_{2}/\sqrt{2},\sigma_{3}/\sqrt{2}\}$. The channel is in
    \emph{diagonal Lorentz normal form} when it is unital and every
    off-diagonal entry of $\widehat{T'}$ is zero.
\end{definition}

\begin{definition}[Non-diagonal Lorentz normal form]\label{def:lorentz_non_diagonal}
    \lean{Wolf.IsLorentzNonDiagonal}\leanok
    For a completely positive, trace-preserving qubit channel
    $T' : \MN{2} \to \MN{2}$, write $\widehat{T'}$ in the normalized
    Pauli basis. The channel is in \emph{non-diagonal Lorentz normal form}
    when, for some $x \in [0,1]$,
    \begin{align}
        \widehat{T'}
        &=
        \begin{pmatrix}
            1   & 0          & 0          & 0   \\
            0   & x/\sqrt{3} & 0          & 0   \\
            0   & 0          & x/\sqrt{3} & 0   \\
            2/3 & 0          & 0          & 1/3
        \end{pmatrix}.
        \label{eq:representations_lorentz_nondiagonal}
    \end{align}
    Trace preservation supplies the first row of the displayed matrix.
\end{definition}

\begin{definition}[Singular Lorentz normal form]\label{def:lorentz_singular}
    \lean{Wolf.IsLorentzSingular}\leanok
    For a completely positive, trace-preserving qubit channel
    $T' : \MN{2} \to \MN{2}$, write $\widehat{T'}$ in the normalized
    Pauli basis. The channel is in \emph{singular Lorentz normal form} when
    \begin{align}
        \widehat{T'}
        &=
        \begin{pmatrix}
            1 & 0 & 0 & 0 \\
            0 & 0 & 0 & 0 \\
            0 & 0 & 0 & 0 \\
            1 & 0 & 0 & 0
        \end{pmatrix},
        \label{eq:representations_lorentz_singular}
    \end{align}
    equivalently, every input state is mapped to $(1+\sigma_{3})/2$.
    Trace preservation supplies the first row of the displayed matrix.
\end{definition}

\begin{theorem}[Displayed non-diagonal and singular representatives]
    \label{thm:lorentz_displayed_qubit_representatives}
    \uses{def:lorentz_non_diagonal, def:lorentz_singular}
    \lean{Wolf.nonDiagonalKrausBase,
        Wolf.nonDiagonalKrausCoefficient,
        Wolf.nonDiagonalKraus,
        Wolf.nonDiagonalMap,
        Wolf.nonDiagonalKraus_isTP,
        Wolf.nonDiagonalMap_isChannel,
        Wolf.nonDiagonalMap_apply,
        Wolf.pauliTransferMatrix_nonDiagonalMap,
        Wolf.isLorentzNonDiagonal_nonDiagonalMap,
        Wolf.choiRank_nonDiagonalMap_eq_three,
        Wolf.nonDiagonalBoundaryKraus,
        Wolf.nonDiagonalKraus_one_one,
        Wolf.nonDiagonalMap_one_eq_boundaryMap,
        Wolf.choiRank_nonDiagonalMap_one,
        Wolf.singularKraus,
        Wolf.singularMap,
        Wolf.singularKraus_isTP,
        Wolf.singularMap_isChannel,
        Wolf.singularMap_apply,
        Wolf.pauliTransferMatrix_singularMap,
        Wolf.isLorentzSingular_singularMap,
        Wolf.choiRank_singularMap}
    \leanok
    For every $x\in[0,1]$, the three Kraus operators displayed in
    \cite[Theorem~8, Eq.~(19)]{VerstraeteVerschelde2002QuantumChannels}
    define the non-diagonal representative
    \begin{align}
        \Delta
        &=\operatorname{diag}(x/\sqrt{3},x/\sqrt{3},1/3),
        &v&=(0,0,2/3).
    \end{align}
    Its Choi rank, equivalently its minimal Kraus rank, is $3$ when
    $x<1$ and $2$ when $x=1$.  The singular representative is the
    trace-to-state channel
    \begin{align}
        X\longmapsto \tr(X)\,\lvert0\rangle\!\langle0\rvert;
    \end{align}
    it has $\Delta=0$, $v=(0,0,1)$, and Choi/Kraus rank $2$.

    These are constructions and rank calculations for the representatives
    in \cite[Theorem~8, Eqs.~(17)--(19)]{VerstraeteVerschelde2002QuantumChannels}
    and Wolf Proposition~2.11, not the Lorentz-orbit classification or the
    derivation that the non-diagonal parameter must lie in $[0,1]$.
    The normalized Choi convention is
    $\tau=\frac14\sum_{ij}\widehat T_{ij}\sigma_i\otimes\sigma_j^{\mathsf T}$;
    hence a raw Pauli correlation matrix has a sign change in the
    $\sigma_2$ input column and is not silently identified with
    $\widehat T$.
\end{theorem}

% Pending Lean formalization: there is currently no declaration for this theorem.
\begin{theorem}[Lorentz normal form for qubit channels]\label{thm:lorentz_normal_form_qubit}
    \uses{thm:exists_normal_form_generic,
        def:invertible_filtering, def:lorentz_diagonal,
        def:lorentz_non_diagonal, def:lorentz_singular,
        thm:lorentz_displayed_qubit_representatives}
    For every qubit channel $T : \MN{2} \to \MN{2}$, there exist invertible
    completely positive maps $\Phi_{1},\Phi_{2}$, both of Kraus rank one, such
    that the filtered channel $T'=\Phi_{2}\circ T\circ\Phi_{1}$ is in one
    of the three Lorentz normal forms: diagonal, non-diagonal, or singular.
    These general filters include scalar freedom and are not restricted to
    determinant-one $\SL(2,\C)$ filterings. Indeed, writing $X_i=c_iS_i$
    with $S_i\in\SL(2,\C)$ separates the Lorentz action from the positive
    map scalar $|c_1c_2|^2$; only the latter can normalize the resulting
    representative to a channel. A proof still requires this scalar
    normalization and the classification of Lorentz orbits.  The displayed
    non-diagonal and singular representatives and their ranks are already
    supplied by Theorem~\ref{thm:lorentz_displayed_qubit_representatives}.
\end{theorem}

\begin{definition}[Pauli--Minkowski identification]
    \label{def:pauli_minkowski_identification}
    \lean{Wolf.pauliMatrices_zero, Wolf.pauliMatrices_succ,
        Wolf.pauliMatrices_isHermitian,
        Wolf.trace_pauliMatrices_mul_pauliMatrices,
        Wolf.pauliMatrixOfMinkowski,
        Wolf.pauliMatrixOfMinkowski_isHermitian,
        Wolf.minkowskiQuadratic, Wolf.minkowskiBilinear,
        Wolf.minkowskiMetric,
        Wolf.minkowskiQuadratic_eq_bilinear_self,
        Wolf.minkowskiBilinear_eq_dotProduct_metric_mulVec,
        Wolf.minkowskiBilinear_polarization,
        Wolf.minkowskiMetric_mul_self,
        Wolf.trace_pauliMatrixOfMinkowski, Wolf.InFutureCone,
        Wolf.trace_mul_eq_ofReal_re_of_isHermitian,
        Wolf.pauliMinkowskiCoordinate,
        Wolf.pauliMinkowskiCoordinate_pauliMatrixOfMinkowski,
        Wolf.coe_pauliMinkowskiCoordinate,
        Wolf.pauliMatrixOfMinkowski_pauliMinkowskiCoordinate,
        Wolf.pauliMatrixOfMinkowskiLinearMap,
        Wolf.pauliMinkowskiCoordinateLinearMap,
        Wolf.pauliMinkowskiEquiv,
        Wolf.pauliMinkowskiEquiv_apply,
        Wolf.pauliMinkowskiEquiv_symm_apply}
    \leanok
    Every Hermitian matrix $M\in M_2^\dagger(\C)$ has unique real Pauli
    coordinates $x=(x_0,x_1,x_2,x_3)$ such that
    \begin{align}
        M&=M(x):=\sum_{i=0}^{3}x_i\sigma_i,
        &x_i&=\frac12\tr(\sigma_iM).
        \label{eq:representations_pauli_minkowski_coordinates}
    \end{align}
    This identifies $M_2^\dagger(\C)$ with $\R^4$.  Write
    \begin{align}
        q(x)&=x_0^2-x_1^2-x_2^2-x_3^2,
        &\eta&=\diag(1,-1,-1,-1),
        \label{eq:representations_minkowski_form}\\
        \mathcal C_+&=\{x\in\R^4:x_0\geq0,\ q(x)\geq0\}
        \label{eq:representations_future_cone}
    \end{align}
    for the Minkowski quadratic form, its matrix, and the closed future cone.
    These are the coordinates used immediately before
    \cite[Eq.~(2.41)]{Wolf2012Quantum}; see also the local source
    \path{Notes/WolfNoteTexSource/ch02_representations.tex}, lines~1040--1044.
\end{definition}

\begin{theorem}[Determinant form and the future cone]
    \label{thm:pauli_minkowski_determinant}
    \lean{Wolf.det_pauliMatrixOfMinkowski,
        Wolf.posSemidef_pauliMatrixOfMinkowski_iff}
    \leanok
    \uses{def:pauli_minkowski_identification}
    For every $x\in\R^4$,
    \begin{align}
        \det M(x)&=q(x)=\langle x,\eta x\rangle,
        \label{eq:representations_pauli_minkowski_determinant}\\
        M(x)\geq0&\iff x\in\mathcal C_+.
        \label{eq:representations_pauli_future_cone}
    \end{align}
\end{theorem}

\begin{proof}\leanok
    The first identity follows by substituting the four Pauli matrices into
    the determinant of a $2\times2$ matrix. For the second, a positive
    semidefinite Hermitian matrix has nonnegative trace and determinant, which
    give $x_0\geq0$ and $q(x)\geq0$. Conversely, the two real eigenvalues of
    $M(x)$ have sum $2x_0\geq0$ and product $q(x)\geq0$, so both are
    nonnegative.
\end{proof}

\begin{definition}[Spinor action on Minkowski space]
    \label{def:spinor_action_minkowski}
    \lean{Wolf.sl2CongruenceLinearEquiv, Wolf.spinorLinearEquiv,
        Wolf.spinorMatrix, Wolf.spinorMatrix_mulVec,
        Wolf.pauliMatrixOfMinkowski_spinorMatrix_mulVec,
        Wolf.pauliMatrixOfMinkowski_single, Wolf.spinorMatrix_apply,
        Wolf.IsSpecialOrthochronousLorentz}
    \leanok
    \uses{def:pauli_minkowski_identification}
    For $X\in\SL(2,\C)$, let $L(X)$ be the real linear transformation of
    $\R^4$ determined by Wolf's Hermitian congruence
    \begin{align}
        M(x)&\longmapsto XM(x)X^\dagger=M(L(X)x).
        \label{eq:representations_spinor_congruence}
    \end{align}
    Its Pauli-basis entries are
    \begin{align}
        L(X)_{ij}
        &=\frac12\tr\!\left(\sigma_iX\sigma_jX^\dagger\right).
        \label{eq:representations_spinor_entries}
    \end{align}
    The trace is real.  This is the four-dimensional action in
    \cite[Eq.~(2.41)]{Wolf2012Quantum}; it is not the three-dimensional
    adjoint action $M\mapsto UMU^{-1}$ on traceless Pauli matrices.
    The special orthochronous Lorentz group is described, in the row-action
    convention of the source, by
    \begin{align}
        \mathrm{SO}^+(1,3)
        &=\{L\in M_4(\R):\det L=1,\ L\eta L^{\mathsf T}=\eta,\ L_{00}>0\}.
        \label{eq:representations_special_orthochronous_lorentz}
    \end{align}
\end{definition}

\begin{theorem}[The spinor image is special orthochronous Lorentz]
    \label{thm:spinor_image_special_orthochronous}
    \lean{Wolf.spinorMatrix_zero_zero,
        Wolf.spinorMatrix_zero_zero_pos,
        Wolf.spinorMatrix_preserves_minkowskiQuadratic,
        Wolf.spinorMatrix_preserves_minkowskiBilinear,
        Wolf.spinorMatrix_transpose_mul_metric_mul,
        Wolf.spinorMatrix_mem_futureCone_iff,
        Wolf.spinorLinearEquiv_mul_apply, Wolf.spinorMatrix_mul,
        Wolf.spinorMatrix_one, Wolf.spinorLinearEquivMap,
        Wolf.spinorMatrix_det, Wolf.spinorMatrix_neg,
        Wolf.spinorMatrix_mul_metric_mul_transpose,
        Wolf.spinorMatrix_isSpecialOrthochronousLorentz,
        Wolf.spinorMap}
    \leanok
    \uses{thm:pauli_minkowski_determinant,
        def:spinor_action_minkowski}
    The assignment $X\mapsto L(X)$ is multiplicative, and for every
    $X\in\SL(2,\C)$,
    \begin{align}
        \det L(X)&=1,
        &L(X)\eta L(X)^{\mathsf T}&=\eta,
        &L(X)_{00}&>0.
        \label{eq:representations_spinor_lorentz_membership}
    \end{align}
    Hence $L(X)\in\mathrm{SO}^+(1,3)$.  It preserves both the Minkowski
    form and the closed future cone, and $L(-X)=L(X)$.
\end{theorem}

\begin{proof}\leanok
    \uses{thm:pauli_minkowski_determinant}
    Congruence by $X$ preserves the determinant because $\det X=1$;
    polarization then gives preservation of the Minkowski bilinear form.
    It preserves the future cone because congruence by an invertible matrix
    preserves positive semidefiniteness.  Moreover,
    $L(X)_{00}=\tr(XX^\dagger)/2>0$.  The determinant of $L(X)$ is a real
    multiplicative character of $\SL(2,\C)$ and is therefore trivial because
    $\SL(2,\C)$ equals its commutator subgroup.  Associativity of congruence
    gives multiplicativity, and the two scalar signs cancel in
    $(-X)M(-X)^\dagger$.
\end{proof}

\begin{definition}[Spinor covering homomorphism]
    \label{def:spinor_cover_homomorphism}
    \lean{Wolf.specialOrthochronousLorentzGroup, Wolf.spinorCoverHom}
    \leanok
    \uses{def:spinor_action_minkowski}
    Regard $\mathrm{SO}^+(1,3)$ from
    (\ref{eq:representations_special_orthochronous_lorentz}) as a matrix
    group.  The spinor action defines the homomorphism
    \begin{align}
        \Lambda:\SL(2,\C)&\longrightarrow\mathrm{SO}^+(1,3),
        &\Lambda(X)&=L(X).
        \label{eq:representations_spinor_cover_hom}
    \end{align}
\end{definition}

\begin{theorem}[Spinor epimorphism and two-point fibres]
    \label{thm:spinor_epimorphism_two_point_fibres}
    \lean{Wolf.IsSpecialOrthochronousLorentz.inv,
        Wolf.IsSpecialOrthochronousLorentz.mul,
        Wolf.row_norm_of_lorentz,
        Wolf.col_norm_of_lorentz,
        Wolf.lorentzRotationBlock_one,
        Wolf.exists_so3_block_of_fixes_time,
        Wolf.boostSpinor_isHermitian,
        Wolf.boostSpinor_posDef,
        Wolf.exists_sl2_spinorMatrix_eq,
        Wolf.spinorCoverHom_surjective,
        Wolf.spinorMatrix_eq_one_iff,
        Wolf.spinorCoverHom_eq_iff_eq_or_eq_neg,
        Wolf.spinorCoverHom_eq_one_iff}
    \leanok
    \uses{def:spinor_cover_homomorphism}
    The homomorphism $\Lambda$ in
    (\ref{eq:representations_spinor_cover_hom}) is surjective.  Its fibres
    contain exactly two points: for all $X,Y\in\SL(2,\C)$,
    \begin{align}
        \Lambda(X)=\Lambda(Y)
        &\iff X=Y\ \text{or}\ X=-Y,
        &\ker\Lambda&=\{I,-I\}.
        \label{eq:representations_spinor_two_point_fibres}
    \end{align}
    Thus $\SL(2,\C)$ is a double cover of $\mathrm{SO}^+(1,3)$, as stated
    after \cite[Eq.~(2.42)]{Wolf2012Quantum}.
\end{theorem}

\begin{proof}\leanok
    \uses{thm:spinor_image_special_orthochronous, thm:pauli_minkowski_determinant}
    Given $L\in\mathrm{SO}^+(1,3)$, set $u=Le_0$.  The Lorentz identities
    give $q(u)=1$ and $u_0>0$.  The canonical boost $B_u$, whose first column
    is $u$, has the positive-definite spinor lift
    $P_u=(M(u)+I)/\sqrt{2(1+u_0)}>0$.  This is the $P>0$ condition in Wolf's
    polar decomposition immediately before \cite[Eq.~(2.44)]{Wolf2012Quantum}.
    The matrix $B_u^{-1}L$ fixes $e_0$;
    its Lorentz equations force the block form $1\oplus R$ with
    $R\in\mathrm{SO}(3)$.  Lift $R$ to $U\in\mathrm{SU}(2)$.  Multiplicativity
    then gives $\Lambda(P_uU)=B_u(1\oplus R)=L$.

    If $\Lambda(X)=I$, congruence applied to $M=I$ gives $XX^\dagger=I$.
    Congruence applied to the Pauli basis then makes $X$ commute with every
    Pauli matrix, so $X$ is scalar.  The equation $\det X=1$ gives
    $X=\pm I$.  Applying this kernel computation to $XY^{-1}$ proves
    (\ref{eq:representations_spinor_two_point_fibres}).
\end{proof}

\begin{theorem}[Boost and rotation spinor exponentials]
    \label{thm:spinor_boost_rotation_exponentials}
    \lean{Wolf.pauliVector, Wolf.boostGenerator, Wolf.rotationGenerator,
        Wolf.exp_smul_pauliVector, Wolf.exp_smul_boostGenerator,
        Wolf.exp_neg_I_smul_pauliVector,
        Wolf.exp_smul_rotationGenerator, Wolf.rapidityMinkowski,
        Wolf.lorentzBoost_rapidityMinkowski_eq_exp,
        Wolf.boostExpSL2, Wolf.rotationExpSL2,
        Wolf.boostExpSL2_coe_eq_exp, Wolf.rotationExpSL2_coe_eq_exp,
        Wolf.spinorMatrix_boostExpSL2,
        Wolf.spinorMatrix_rotationExpSL2}
    \leanok
    \uses{def:spinor_cover_homomorphism}
    Let $n\in\R^3$ be a unit vector.  Write $n\cdot\sigma$ for its Pauli
    contraction, $n\cdot B$ for the boost generator with time--space blocks
    $n$ and $n^{\mathsf T}$, and $n\cdot R$ for Wolf's rotation generator
    $R_i=\sum_{j,k}\varepsilon_{ijk}\lvert k\rangle\langle j\rvert$.
    For every $t\in\R$,
    \begin{align}
        \begin{aligned}
            \exp(t\,n\cdot\sigma)
            &=\cosh(t)I+\sinh(t)n\cdot\sigma,\\
            \exp(t\,n\cdot B)
            &=I+(\cosh(t)-1)(n\cdot B)^2+\sinh(t)n\cdot B,\\
            \exp(-it\,n\cdot\sigma)
            &=\cos(t)I-i\sin(t)n\cdot\sigma,\\
            \exp(t\,n\cdot R)
            &=I+(1-\cos(t))(n\cdot R)^2+\sin(t)n\cdot R.
        \end{aligned}
        \label{eq:representations_spinor_exponential_closed_forms}
    \end{align}
    If $u(n,t)=(\cosh(t),\sinh(t)n)$, then
    $B_{u(n,t)}=\exp(t\,n\cdot B)$, and the spinor map satisfies
    \begin{align}
        \begin{aligned}
            \Lambda\!\left(\exp\!\left(\frac{t}{2}n\cdot\sigma\right)\right)
            &=\exp(t\,n\cdot B),\\
            \Lambda\!\left(\exp\!\left(-\frac{it}{2}n\cdot\sigma\right)\right)
            &=\exp(t\,n\cdot R).
        \end{aligned}
        \label{eq:representations_spinor_exponential_map}
    \end{align}
    The positive sign in the second Lorentz exponential is the correction to
    the sign printed in \cite[Eq.~(2.44)]{Wolf2012Quantum}; with the displayed
    Pauli matrices and column-vector convention, the printed negative sign
    rotates in the opposite direction.
\end{theorem}

\begin{proof}\leanok
    The Pauli contraction satisfies $(n\cdot\sigma)^2=I$, the boost generator
    satisfies $(n\cdot B)^3=n\cdot B$, and the rotation generator satisfies
    $(n\cdot R)^3=-n\cdot R$.  Splitting each exponential series into its even
    and odd terms gives the four closed forms in
    (\ref{eq:representations_spinor_exponential_closed_forms}).  Direct
    substitution of $u(n,t)$ into the canonical boost gives
    $B_{u(n,t)}=\exp(t\,n\cdot B)$.  The two half-parameter matrices have
    determinant one, and evaluating their Pauli congruence entries with the
    double-angle identities gives
    (\ref{eq:representations_spinor_exponential_map}).
\end{proof}

\begin{theorem}[Pauli transfer matrix under two-sided filtering]
    \label{thm:pauli_transfer_two_sided_filtering}
    \lean{Wolf.spinorMatrixComplex, Wolf.pauli_expansion_four,
        Wolf.coe_spinorMatrix_apply,
        Wolf.sl2Congruence_pauli,
        Wolf.pauliTransferEntry_unitaryConj_comp,
        Wolf.pauliTransferEntry_comp_unitaryConj,
        Wolf.pauliTransferMatrix_unitaryConj_comp,
        Wolf.pauliTransferMatrix_comp_unitaryConj,
        Wolf.pauliTransferMatrix_two_sided_filtering,
        Wolf.coe_pauliTransferMatrixReal_of_preservesHermiticity,
        Wolf.pauliTransferMatrixReal_two_sided_filtering,
        Wolf.SLFiltering.toSL2,
        Wolf.pauliTransferMatrixReal_slFiltering}
    \leanok
    \uses{def:pauli_transfer_matrix, def:spinor_action_minkowski}
    Let $T:M_2(\C)\to M_2(\C)$ be complex linear and let
    $X_1,X_2\in\SL(2,\C)$. If $L_i=L(X_i)$ and $L_{i,\C}$ denotes its
    complex scalar extension, then
    \begin{align}
        \widehat{\Phi_{X_2}\circ T\circ\Phi_{X_1}}
        &=L_{2,\C}\,\widehat T\,L_{1,\C}.
        \label{eq:representations_pauli_transfer_spinor_action}
    \end{align}
    If $T$ preserves Hermiticity, then $\widehat T$ has real entries, and
    Equation~\ref{eq:representations_pauli_transfer_spinor_action} is the real
    identity
    \begin{align}
        \widehat{\Phi_{X_2}\circ T\circ\Phi_{X_1}}
        &=L_2\,\widehat T\,L_1,
        \label{eq:representations_pauli_transfer_spinor_action_real}
    \end{align}
    in the exact order of \cite[Eq.~(2.43)]{Wolf2012Quantum}.
\end{theorem}

\begin{proof}\leanok
    Expand each matrix in the four-Pauli basis using the trace pairing in
    Equation~\ref{eq:pauli_transfer_entries}. Postfiltering by $\Phi_{X_2}$
    multiplies each transfer-matrix column on the left by $L_{2,\C}$, while
    prefiltering by $\Phi_{X_1}$ multiplies each row on the right by
    $L_{1,\C}$. Combining these two finite sums gives
    Equation~\ref{eq:representations_pauli_transfer_spinor_action}. For a
    Hermiticity-preserving map, the trace of the product of the two Hermitian
    matrices $\sigma_i$ and $T(\sigma_j)$ is real, giving the stated real
    specialization.
\end{proof}

%% The following sections will not be needed for the fundamental theorem of the matrix product states.

%% =========================================================
\section{Determinant of a quantum channel}
\label{sec:channel_determinant}
%% =========================================================

%% Source: \cite[Section~6.1, Theorem~6.1(2)]{Wolf2012Quantum}.

\begin{definition}[Channel determinant]\label{def:channel_det}
    \lean{channelDet}\leanok
    For a linear map $T : \MN{D} \to \MN{D}$, the
    \emph{channel determinant} $\det T$ is the determinant of the
    matrix of $T$ with respect to the standard matrix-unit basis
    of $\MN{D}$.
    This is the quantity studied in \cite[Section~6.1]{Wolf2012Quantum}.
\end{definition}

\begin{remark}[Channel determinant as an operator determinant]
    \label{rem:channel_det_linearMap_det}
    \lean{channelDet_eq_linearMap_det}
    \leanok
    The channel determinant agrees with the ordinary determinant of the linear
    endomorphism $T : \MN{D} \to \MN{D}$.
\end{remark}

\begin{definition}[Unitary channel]\label{def:unitary_channel}
    \lean{unitaryChannel, unitaryChannel_isChannel}\leanok
    For a unitary matrix $U \in \mc{U}(D)$, the
    \emph{unitary channel} is $T(\rho)=U\rho U^\dagger$.
    It is automatically a quantum channel.
\end{definition}

\begin{theorem}[Unitary channels have determinant one]
    \label{thm:channel_det_unitary_eq_one}
    \lean{channelDet_unitary_eq_one, channelDet_norm_eq_one_of_unitaryChannel}
    \leanok
    \uses{def:channel_det, def:unitary_channel}
    If $T(\rho)=U\rho U^\dagger$ is a unitary channel, then
    $\det T=1$ and hence $|\det T|=1$.
\end{theorem}

\begin{proof}\leanok
    Vectorization identifies $T$ with a Kronecker product
    $\overline{U}\otimes U$, whose determinant is
    $\overline{\det U}^{\,D}(\det U)^D=1$.
\end{proof}

\begin{theorem}[Eigenvalue disk bound for positive trace-preserving maps]
    \label{thm:positive_tp_eigenvalue_disk}
    \lean{IsPositiveMap.eigenvalue_norm_le_one_of_tracePreserving}\leanok
    \uses{def:positive_map, def:trace_preserving}
    Let $D>0$, and let $T:\MN{D}\to\MN{D}$ be positive and trace-preserving.
    If $T(X)=\mu X$ for some nonzero $X$, then $|\mu|\leq1$.
    This is the unit-disk conclusion of
    \cite[Proposition~6.1]{Wolf2012Quantum}.
\end{theorem}

\begin{proof}\leanok
    If $\tr(X)\neq0$, trace preservation gives $\mu=1$. Otherwise decompose
    $X$ into trace-zero Hermitian and skew-Hermitian parts. Each Hermitian
    part is a scalar multiple of the difference of two density matrices.
    Positivity and trace preservation keep all iterates of those density
    matrices in the compact set of density matrices, so the orbit of $X$ is
    bounded. Since $T^n(X)=\mu^nX$, this is impossible when $|\mu|>1$.
\end{proof}

\begin{theorem}[Real determinant interval for positive trace-preserving maps]
    \label{thm:channelDet_norm_le_one}
    \lean{channelDet_norm_le_one_of_positive_tracePreserving,
        channelDet_norm_le_one_of_channel,
        ChannelDeterminant.Internal.channelDet_star_eq_of_map_conjTranspose,
        ChannelDeterminant.Internal.channelDet_star_eq_of_isPositiveMap,
        ChannelDeterminant.Internal.exists_real_channelDet_mem_Icc_of_positive_tracePreserving}\leanok
    \uses{def:channel_det, def:positive_map, def:trace_preserving}
    For any positive trace-preserving map $T : \MN{D} \to \MN{D}$,
    $\det T$ is real and belongs to $[-1,1]$; in particular,
    $|\det T|\leq1$.
    This is \cite[Theorem~6.1(1)]{Wolf2012Quantum}.
\end{theorem}

\begin{proof}\leanok
    \uses{thm:positive_tp_eigenvalue_disk}
    Positivity implies Hermiticity preservation.  In the matrix-unit basis,
    conjugating the channel matrix entrywise is the same as conjugating the
    map by the pair-swap involution induced by $X\mapsto X^{\mathsf T}$.
    Determinant multiplicativity and involutivity therefore give
    $\overline{\det T}=\det T$, which is the source's conjugate-pair reality
    argument in coordinates.

    Every eigenvalue $\mu$ of $T$ satisfies $|\mu|\leq1$
    (positivity and trace preservation force spectral radius~$\leq1$).
    Since $\det T$ is the product of the eigenvalues
    (counted with algebraic multiplicity), the bound
    $|\det T|=\prod_i|\mu_i|\leq1$ follows by induction.
\end{proof}

\begin{theorem}[Determinant saturation for positive trace-preserving maps]
    \label{thm:channelDet_norm_eq_one_iff_positive_tp}
    \lean{ChannelDeterminant.Internal.channel_all_eigenvalues_norm_one_of_positive_tracePreserving,
        ChannelDeterminant.Internal.channel_all_eigenvalues_norm_one,
        Module.End.peripheralSubspace_eq_top_of_all_eigenvalues_norm_one,
        Module.End.peripheralProjection_eq_one_of_all_eigenvalues_norm_one,
        ChannelDeterminant.Internal.mapsPSDConeOnto_of_channelDet_norm_eq_one,
        ChannelDeterminant.Internal.exists_unitary_or_transpose_of_channelDet_norm_eq_one,
        ChannelDeterminant.Internal.channelDet_transpose_norm_eq_one,
        ChannelDeterminant.Internal.channelDet_norm_eq_one_of_unitaryChannel_comp_transpose,
        ChannelDeterminant.Internal.channelDet_norm_eq_one_iff_exists_unitary_or_transpose_of_positive_tracePreserving}\leanok
    \uses{def:channel_det, def:positive_map, def:trace_preserving,
        def:unitary_channel}
    Let $D>0$ and let $T:\MN{D}\to\MN{D}$ be positive and
    trace-preserving.  Then
    \begin{align}
        |\det T|=1
        \quad\Longleftrightarrow\quad
        \exists U\in\mc{U}(D),\quad
        T(A)=UAU^\dagger
        \ \text{or}\
        T(A)=UA^{\mathsf T}U^\dagger .
    \end{align}
    This is \cite[Theorem~6.1(2)]{Wolf2012Quantum}.
\end{theorem}

\begin{proof}\leanok
    \uses{thm:positive_tp_eigenvalue_disk,
        thm:peripheral_projection_recurrent_subsequence, thm:wolf_prop_3_6,
        thm:channel_det_unitary_eq_one}
    Saturation forces every eigenvalue to have modulus one, so the peripheral
    subspace is all of $\MN{D}$ and the peripheral projection $T_\phi$ is the
    identity.  Along Wolf's Dirichlet recurrent subsequence,
    $T^{n_i}\to T_\phi=\id$.  Since $T$ is bijective, the preimage $X$ of a
    positive matrix $A$ is the limit of the positive matrices
    $T^{n_i-1}(A)$; closedness of the positive-semidefinite cone makes
    $T^{-1}$ positive.

    Thus $T$ maps the positive-semidefinite cone onto itself.  The already
    formalized cone/rank/Wigner route of Wolf Proposition~3.6 gives
    $T(A)=YAY^\dagger$ or $T(A)=YA^{\mathsf T}Y^\dagger$ with $Y$
    invertible.  Trace preservation and nondegeneracy of the trace pairing
    imply $Y^\dagger Y=\Id$, so $Y$ is unitary.  Conversely, unitary
    conjugation has determinant one, while transposition is an involution and
    hence has determinant of modulus one.
\end{proof}

\begin{theorem}[Sign of transposition and the extremal negative determinant]
    \label{thm:channelDet_transpose_sign}
    \lean{ChannelDeterminant.Internal.matrixBasisTransposePerm,
        ChannelDeterminant.Internal.matrixBasisTransposePerm_apply,
        ChannelDeterminant.Internal.channelMatrix_transposeLinearMapComplex,
        ChannelDeterminant.Internal.channelDet_transposeLinearMapComplex,
        ChannelDeterminant.Internal.odd_mul_pred_div_two_iff_odd_div_two,
        ChannelDeterminant.Internal.channelDet_transposeLinearMapComplex_eq_neg_one_iff,
        ChannelDeterminant.Internal.channelDet_eq_neg_one_iff_exists_unitary_transpose_of_positive_tracePreserving,
        ChannelDeterminant.Internal.channelDet_eq_one_iff_exists_unitary_of_positive_tracePreserving_of_odd}\leanok
    \uses{thm:channelDet_norm_eq_one_iff_positive_tp,
        thm:channel_det_unitary_eq_one}
    Ordinary transposition on $\MN{D}$ satisfies
    \begin{align}
        \det(A\mapsto A^{\mathsf T})
        =(-1)^{D(D-1)/2}.
    \end{align}
    Consequently it has determinant $-1$ exactly when
    $\lfloor D/2\rfloor$ is odd.  For a positive trace-preserving $T$ in
    positive dimension, $\det T=-1$ exactly when that parity condition holds
    and $T(A)=UA^{\mathsf T}U^\dagger$ for a unitary $U$.  In those odd-parity
    dimensions, $\det T=1$ exactly for unitary conjugations.  This is
    \cite[Theorem~6.1(3)]{Wolf2012Quantum}.
\end{theorem}

\begin{proof}\leanok
    In the matrix-unit basis, transposition is the permutation that exchanges
    $E_{ij}$ and $E_{ji}$.  It fixes the $D$ diagonal units and partitions the
    remaining $D^2-D$ units into $D(D-1)/2$ transpositions, giving the displayed
    sign.  Splitting $D$ into its even and odd cases proves that
    $D(D-1)/2$ and $\lfloor D/2\rfloor$ have the same parity.  The source-facing
    classifications then follow from determinant saturation and the fact that
    every unitary conjugation has determinant one.
\end{proof}

\begin{theorem}[Determinant under composition]
    \label{thm:channelDet_comp}
    \lean{channelDet_comp}\leanok
    \uses{def:channel_det}
    For linear maps $T_1,T_2:\MN{D}\to\MN{D}$,
    \begin{align}
        \det(T_1T_2)=\det(T_1)\det(T_2),
    \end{align}
    where $T_1T_2=T_1\circ T_2$, so $T_2$ acts first.  This is
    \cite[Equation~(6.22)]{Wolf2012Quantum}.
\end{theorem}

\begin{proof}\leanok
    Rewrite the channel determinant as the determinant of the underlying
    linear endomorphism and apply multiplicativity of the latter.
\end{proof}

\begin{theorem}[Determinant monotonicity under composition: algebraic part]
    \label{thm:channelDet_comp_monotone_algebraic}
    \lean{channelDet_norm_comp_le_of_positive_tracePreserving,
        channelDet_norm_comp_eq_iff}\leanok
    \uses{def:channel_det, def:positive_map, def:trace_preserving}
    Let $T_1,T_2:\MN{D}\to\MN{D}$ be positive and trace-preserving.  Then
    \begin{align}
        |\det(T_1T_2)|\leq |\det T_1|.
    \end{align}
    For arbitrary linear maps, the equality
    $|\det(T_1T_2)|=|\det T_1|$ holds exactly when
    \begin{align}
        \det T_1=0 \quad\text{or}\quad |\det T_2|=1.
    \end{align}
    This is the multiplicative and determinant-bound part of the
    monotonicity corollary following
    \cite[Theorem~6.1]{Wolf2012Quantum}.
\end{theorem}

\begin{proof}\leanok
    \uses{thm:channelDet_comp, thm:channelDet_norm_le_one}
    Theorem~\ref{thm:channelDet_comp} gives
    $|\det(T_1T_2)|=|\det T_1|\,|\det T_2|$.  The determinant bound for
    $T_2$ proves the inequality.  Cancelling the nonzero factor
    $|\det T_1|$ proves the equality split; when that factor is zero,
    equality is automatic.

    The source's further replacement of $|\det T_2|=1$ by ``unitary
    conjugation or matrix transposition'' is exactly
    Theorem~\ref{thm:channelDet_norm_eq_one_iff_positive_tp}; the present
    declaration isolates the algebraic equality split so the saturation
    classification remains reusable.
\end{proof}

\begin{definition}[Inverse of a bijective matrix endomorphism]
    \label{def:inverse_of_bijective_matrix_endomorphism}
    \lean{ChannelDeterminant.Internal.inverseOfBijective,
        ChannelDeterminant.Internal.apply_inverseOfBijective,
        ChannelDeterminant.Internal.inverseOfBijective_apply,
        ChannelDeterminant.Internal.comp_inverseOfBijective,
        ChannelDeterminant.Internal.inverseOfBijective_comp}
    \leanok
    Let $T:\MN{D}\to\MN{D}$ be a bijective complex-linear map.  Write
    $T^{-1}$ for the inverse linear map supplied by the associated linear
    equivalence; it is both a left and a right inverse of $T$.
\end{definition}

\begin{theorem}[Trace preservation and positivity of the inverse]
    \label{thm:positive_inverse_preliminaries}
    \lean{ChannelDeterminant.Internal.inverseOfBijective_isTracePreservingMap,
        ChannelDeterminant.Internal.mapsPSDConeOnto_of_inverseOfBijective_isPositiveMap,
        ChannelDeterminant.Internal.inverseOfBijective_isPositiveMap_of_mapsPSDConeOnto,
        ChannelDeterminant.Internal.inverseOfBijective_isPositiveMap_iff_mapsPSDConeOnto,
        ChannelDeterminant.Internal.channelDet_norm_eq_one_of_inverseOfBijective_isPositiveMap}
    \leanok
    \uses{def:inverse_of_bijective_matrix_endomorphism,
        def:positive_map, def:trace_preserving,
        def:maps_psd_cone_onto, def:channel_det}
    The inverse of a bijective trace-preserving linear map is trace preserving.
    If $T$ is positive, then $T^{-1}$ is positive exactly when $T$ maps the
    positive-semidefinite cone onto itself.  If both maps are positive and
    trace preserving, then $|\det T|=1$.
\end{theorem}

\begin{proof}\leanok
    \uses{thm:channelDet_norm_le_one, thm:channelDet_comp}
    For $X=T^{-1}(Y)$, trace preservation of $T$ gives
    $\tr(T^{-1}(Y))=\tr(Y)$.  Positivity of the inverse is equivalent to
    positivity of the unique preimage of every positive-semidefinite matrix.
    Finally,
    $\det(T)\det(T^{-1})=1$, while the determinant bound places both moduli
    at most one, so both moduli equal one.
\end{proof}

\begin{theorem}[Explicit inverses of the standard forms]
    \label{thm:positive_inverse_standard_forms}
    \lean{ChannelDeterminant.Internal.inverseOfBijective_unitaryChannel,
        ChannelDeterminant.Internal.inverseOfBijective_unitaryChannel_comp_transpose,
        ChannelDeterminant.Internal.inverseOfBijective_unitaryChannel_isPositiveMap,
        ChannelDeterminant.Internal.inverseOfBijective_unitaryChannel_comp_transpose_isPositiveMap}
    \leanok
    \uses{def:inverse_of_bijective_matrix_endomorphism, def:unitary_channel,
        def:positive_map}
    For a unitary $U$, the inverse of $A\mapsto UAU^\dagger$ is conjugation
    by $U^{-1}$.  The inverse of
    $A\mapsto UA^{\mathsf T}U^\dagger$ is ordinary transposition after
    conjugation by $U^{-1}$.  Both inverse maps are positive; the reversed
    composition order in the transpose branch is essential.
\end{theorem}

\begin{proof}\leanok
    \uses{thm:transpose_positive_trace_preserving}
    Compose the displayed maps directly, using involutivity of ordinary
    transposition.  Unitary conjugation, transposition, and their compositions
    preserve the positive-semidefinite cone.
\end{proof}

\begin{theorem}[Positive invertible maps]
    \label{thm:wolf_positive_invertible_maps}
    \lean{ChannelDeterminant.Internal.wolfPositiveInvertibleMaps}
    \leanok
    \uses{def:inverse_of_bijective_matrix_endomorphism,
        def:positive_map, def:trace_preserving, def:unitary_channel}
    Let $D>0$, and let $T:\MN{D}\to\MN{D}$ be positive, trace preserving,
    and bijective.  Then $T^{-1}$ is positive if and only if there is a
    unitary $U$ for which
    \begin{align}
        T(A)=UAU^\dagger
        \qquad\text{or}\qquad
        T(A)=UA^{\mathsf T}U^\dagger.
    \end{align}
    This is the corollary ``Positive invertible maps'' following
    \cite[Theorem~6.1]{Wolf2012Quantum}.
\end{theorem}

\begin{proof}\leanok
    \uses{thm:positive_inverse_preliminaries,
        thm:positive_inverse_standard_forms,
        thm:channelDet_norm_eq_one_iff_positive_tp}
    If $T^{-1}$ is positive, it is automatically trace preserving and the
    determinant bounds for $T$ and $T^{-1}$ force $|\det T|=1$.  Apply the
    source-general determinant-saturation classification.  Conversely, use
    the explicit positive inverse of each standard form above.  This is
    exactly Wolf's determinant route and retains the transpose branch.
\end{proof}

\begin{theorem}[Determinant one iff the channel is unitary]
    \label{thm:channelDet_norm_eq_one_iff_unitary}
    \lean{channelDet_norm_eq_one_iff_exists_unitaryChannel}\leanok
    \uses{def:channel, def:channel_det, def:unitary_channel}
    For a CPTP map $T : \MN{D} \to \MN{D}$,
    \begin{align}
        |\det T|=1
        &\iff
        \exists\,U \in \mc{U}(D),\quad T(\rho)=U\rho U^\dagger.
        \label{eq:representations_det_unitary_iff}
    \end{align}
    This is the CPTP specialization of
    \cite[Theorem~6.1(2)]{Wolf2012Quantum}; complete positivity excludes the
    genuinely transpose-type branch.
\end{theorem}

\begin{proof}
    \leanok
    \uses{thm:channelDet_norm_le_one, thm:channel_det_unitary_eq_one}
    Reverse direction is Theorem~\ref{thm:channel_det_unitary_eq_one}.
    Forward direction: $|\det T|=1$ together with $|\mu|\leq1$ for every
    eigenvalue $\mu$ (from trace preservation and positivity) forces every
    eigenvalue to satisfy $|\mu|=1$. Determinant saturation transfers to the
    unital Heisenberg dual $T^*(Y)=\sum_iK_i^\dagger YK_i$. The
    determinant--Hilbert--Schmidt bound and a trace-summing argument then force
    equality in the Kadison--Schwarz inequality on the standard matrix basis.
    The resulting Kraus commutation relations make $T^*$ multiplicative on all
    matrices, hence a $*$-automorphism of $\MN{D}$. By the Skolem--Noether
    theorem it is inner, and therefore
    $T(\rho)=U\rho U^\dagger$ for some unitary $U$. Kraus freedom then gives
    $K_i=c_iU$ with $\sum_i|c_i|^2=1$.
\end{proof}

\begin{theorem}[Positive determinant for small Kraus rank]
    \label{thm:wolf_positive_det_small_kraus_rank}
    \lean{channelDet_nonneg_of_two_kraus,
        channelDet_nonneg_of_hasKrausRankLE_two,
        IsKrausCP.channelDet_nonneg_of_choiRank_le_two}
    \leanok
    \uses{def:is_kraus_cp, def:kraus_rank_rect, def:channel_det}
    Let $T:\MN{d}\to\MN{d}$ be a completely positive linear map with
    Kraus rank at most two. Then
    \begin{align}
        \det T
        &\geq0.
    \end{align}
    No trace-preservation hypothesis is assumed. This is the proposition
    ``Positive determinant for small Kraus rank'' in
    \cite[Chapter~6, Equation~(6.26)]{Wolf2012Quantum}.
\end{theorem}

\begin{proof}
    \leanok
    \uses{thm:kraus_rect_rank_minimal,
        thm:transfer_matrix_kraus,
        thm:channel_det_transfer_matrix}
    By minimality of Kraus rank and zero-padding, write
    $T(X)=AXA^\dagger+BXB^\dagger$, including Kraus ranks zero and one.
    Wolf writes
    \begin{align}
        \widehat T
        &=A\otimes\overline A+B\otimes\overline B.
    \end{align}
    With the column-stacking convention of
    Definition~\ref{def:mpu_transfer_matrix}, the two
    Kronecker factors are swapped, so
    $\widehat T=\overline A\otimes A+\overline B\otimes B$.

    First suppose $\det A\neq0$ and put $C=A^{-1}B$. Equation~(6.26), in
    this convention, is the exact factorization
    \begin{align}
        \widehat T
        &=(\overline A\otimes A)
          (\Id+\overline C\otimes C).
    \end{align}
    The determinant of the first factor is
    $\overline{(\det A)^d}(\det A)^d\geq0$.
    For the second factor, take a unitary Schur decomposition
    $C=URU^\dagger$ with diagonal entries $\lambda_i=R_{ii}$. Unitary
    similarity reduces its determinant to
    \begin{align}
        \prod_{i,j}
        \bigl(1+\overline{\lambda_i}\lambda_j\bigr).
    \end{align}
    Each off-diagonal factor is paired with its complex conjugate, while the
    diagonal factors are $1+|\lambda_i|^2$, so this product is nonnegative.

    To remove the assumption on $A$ without an unproved density assertion,
    set $\varepsilon_n=1/(n+1)$ and $A_n=A+\varepsilon_n\Id$. If
    $A=URU^\dagger$ is a unitary Schur decomposition, then
    $A_n=U(R+\varepsilon_n\Id)U^\dagger$. Every zero diagonal entry of $R$
    becomes the nonzero number $\varepsilon_n$, while every nonzero diagonal
    entry remains nonzero for all sufficiently large $n$. Hence
    $\det A_n\neq0$ eventually. The nonsingular case applies to $(A_n,B)$;
    since $A_n\to A$ and the determinant of
    $\overline{A_n}\otimes A_n+\overline B\otimes B$ is continuous, passing
    to the limit proves the claim.
\end{proof}
